\documentclass[a4paper,12pt]{article}
\usepackage[T2A]{fontenc}
\usepackage[utf8]{inputenc}
\usepackage[russian]{babel}
\usepackage{amsmath,amssymb,mathtools}
\renewcommand{\familydefault}{\sfdefault}
\usepackage{geometry}
\geometry{left=22mm,right=22mm,top=23mm,bottom=24mm}
\usepackage{microtype}
\usepackage{array,tabularx,booktabs}
\usepackage{enumitem}
\usepackage{parskip}
\usepackage{xcolor}
\usepackage{tikz}
\usepackage{fancyhdr}
\usepackage{hyperref}
\usepackage{setspace}
\definecolor{accent}{HTML}{167E80}
\definecolor{darkaccent}{HTML}{18545C}
\definecolor{textgray}{HTML}{3B454A}
\definecolor{error}{HTML}{B83E42}
\definecolor{linegray}{HTML}{D0D8D8}
\hypersetup{colorlinks=true,linkcolor=darkaccent,pdfauthor={Лёвин Артём Александрович},pdftitle={Ревью домашней работы — Ксения Васильченко, 8 ноября 2026}}
\setlength{\parindent}{0pt}
\setlength{\parskip}{6pt}
\setstretch{1.13}
\setlist[itemize]{leftmargin=*,topsep=2pt,itemsep=3pt}
\pagestyle{fancy}
\fancyhf{}
\fancyhead[L]{\footnotesize\color{textgray} Ревью домашней работы}
\fancyhead[R]{\footnotesize\color{textgray} 08.11.2026}
\fancyfoot[C]{\footnotesize\color{textgray} \thepage}
\renewcommand{\headrulewidth}{0pt}
\renewcommand{\footrulewidth}{0pt}
\newcommand{\tagline}[1]{\par\vspace{0.35em}{\large\bfseries\color{darkaccent} #1}\par\vspace{0.08em}}
\newcommand{\errtag}{\textcolor{error}{\textbf{Ошибка:}}}
\newcommand{\keytag}{\textcolor{darkaccent}{\textbf{Ключевая идея:}}}
\newcommand{\finaltag}{\textcolor{accent}{\textbf{Итог:}}}
\begin{document}
\thispagestyle{empty}
\begin{titlepage}
\centering
\vspace*{38mm}
{\small\color{accent}\bfseries МАТЕМАТИКА}\par
\vspace{13mm}
{\Huge\bfseries\color{darkaccent} Ревью\\[3mm] домашней работы}\par
\vspace{10mm}
{\large Доли, дроби и проценты}\par
\vspace{22mm}
{\Large\bfseries Ксения Васильченко}\par
\vspace{8mm}
{\normalsize 8 ноября 2026 года}\par
\vfill
{\footnotesize\color{textgray} Лёвин Артём Александрович\\[2mm]
эксклюзивно для Ксении Васильченко}\par
\vspace{9mm}
\begin{tikzpicture}
\draw[accent!58,line width=.7pt] (-6,0) .. controls (-2.4,.55) and (1.2,-.4) .. (6,.05);
\end{tikzpicture}
\vspace*{8mm}
\end{titlepage}
\setcounter{page}{2}

{\LARGE\bfseries\color{darkaccent} Сводная таблица}\par
\vspace{2mm}
{\color{textgray} Ниже представлены задания, требующие исправления или завершения решения. Нажмите на номер, чтобы перейти к разбору.}\par
\vspace{9mm}
\renewcommand{\arraystretch}{1.6}
\small
\begin{tabularx}{\linewidth}{@{}>{\centering\arraybackslash}p{10mm}p{27mm}X>{\raggedright\arraybackslash}p{22mm}>{\raggedright\arraybackslash}p{20mm}@{}}
\toprule
\textbf{№} & \textbf{Тема} & \textbf{Суть ошибки} & \textbf{Ваш ответ} & \textbf{Верный ответ}\\
\midrule
\hyperref[task:7]{\textbf{7}} & Дробь от целого & Найдена прочитанная доля книги, но не вычислено число страниц. & Не указан & 60 страниц\\
\addlinespace[2mm]
\hyperref[task:10]{\textbf{10}} & Проценты в задаче на изменение & За основу взяты 30\% вместо оставшихся 70\%; есть ошибка в записи процента дробью. & 210 часов & 90 часов\\
\bottomrule
\end{tabularx}
\normalsize
\vfill
\begin{center}
\textcolor{accent}{\rule{33mm}{0.6pt}}\\[4mm]
{\small\color{textgray} В каждом разборе приведено одно упражнение на закрепление.}
\end{center}
\newpage
\thispagestyle{empty}
\vspace*{0.30\textheight}
\begin{center}
{\Huge\bfseries\color{darkaccent} Разбор ошибок}\par
\vspace{11mm}
{\large Доли целого и процентное изменение}\par
\vspace{14mm}
\textcolor{accent}{\rule{40mm}{0.75pt}}
\end{center}
\newpage
\phantomsection\label{task:7}
{\LARGE\bfseries\color{darkaccent} Задание 7. Сколько страниц в книге?}\par
\vspace{3mm}
\tagline{Условие}
Ученик прочитал половину книги и ещё четверть книги. Всего прочитано 45 страниц. Сколько страниц в книге?

\tagline{Ваше решение}
Вы правильно сложили прочитанные части:
\[
\frac12+\frac14=\frac24+\frac14=\frac34.
\]

\textbf{Ваш ответ:} не указан. После вычисления доли книги решение не завершено.

\tagline{Где остановилось решение}
\errtag\ Получена правильная доля, но не выполнен следующий обязательный шаг: определение целого по известной его части.

\textbf{Последний правильный шаг:} прочитано \(\frac34\) всей книги.

\textbf{Первый пропущенный шаг:} записать, что 45 страниц составляют \(\frac34\) неизвестного количества страниц, и найти это количество.

\tagline{Почему этого недостаточно}
Дробь \(\frac34\) отвечает на вопрос, \emph{какую часть} книги прочитали. В условии требуется узнать, \emph{сколько всего страниц} в книге. Для этого нужно связать долю с конкретным числом страниц.

Если известно, что \(\frac34\) целого равны 45, то сначала можно найти одну четверть, а затем все четыре четверти.

\keytag\ Чтобы найти целое по его дроби, разделите известное значение на эту дробь. Можно также найти значение одной равной части и умножить на число таких частей в целом.

\newpage
\tagline{Исправление от последнего правильного шага}
Три четверти книги составляют 45 страниц. Значит, одна четверть равна
\[
45:3=15\ \text{страниц}.
\]
Вся книга содержит четыре четверти, поэтому
\[
15\cdot4=60\ \text{страниц}.
\]

\tagline{Полное правильное решение}
\begin{enumerate}[leftmargin=*,itemsep=2pt]
\item Приводим дроби к общему знаменателю: \(\frac12=\frac24\).
\item Находим прочитанную долю: \(\frac24+\frac14=\frac34\).
\item Находим всё число страниц: \(45:\frac34=45\cdot\frac43=60\).
\end{enumerate}

\tagline{Проверка}
Половина от 60 страниц равна 30, а четверть от 60 равна 15. Всего прочитано \(30+15=45\) страниц, как в условии.

\finaltag\ В книге \(\boxed{60}\) страниц.

\tagline{Задача на закрепление}
Ученица прочитала треть книги и ещё четверть книги. Всего прочитано 35 страниц. Сколько страниц в книге?

\newpage
\phantomsection\label{task:10}
{\LARGE\bfseries\color{darkaccent} Задание 10. Работа комбайна}\par
\vspace{3mm}
\tagline{Условие}
Новый комбайн выполняет работу за 63 часа, что на 30\% меньше времени, которое требовалось старому комбайну. Сколько часов требовалось старому комбайну?

\tagline{Ваше решение}
В записи встречаются переходы:
\[
30\%=\frac{300}{100}=\frac{3}{10},\qquad
63:\frac{3}{10}=210.
\]
\textbf{Ваш ответ:} 210 часов.

\tagline{Где именно возникает ошибка}
\errtag\ При переводе 30\% в дробь записано \(\frac{300}{100}\) вместо \(\frac{30}{100}\). Кроме того, 63 часа ошибочно приняты за 30\% времени старого комбайна.

\textbf{Последний достоверный исходный факт:} новый комбайн работает 63 часа и тратит на 30\% меньше времени, чем старый.

\textbf{Первый неправильный вычислительный шаг:}
\[
\textcolor{error}{30\%=\frac{300}{100}}.
\]
Действительно, \(\frac{300}{100}=3\), а 30\% равны \(\frac{30}{100}=\frac3{10}\). Даже после исправления этой записи делить 63 на \(\frac3{10}\) нельзя: 63 часа соответствуют оставшейся доле времени.

\tagline{Почему деление на 30\% неверно}
За 100\% принимаем время \emph{старого} комбайна. Новый тратит на 30\% меньше, то есть работает \(100\%-30\%=70\%\) прежнего времени. Поэтому 63 часа равны 70\% старого времени.

\keytag\ Слова «на 30\% меньше» означают, что осталось 70\% исходной величины. Сначала определите, какая величина принята за 100\%, затем установите долю, соответствующую известному числу.

\newpage
\tagline{Исправление}
Правильно переводим процент в дробь и находим оставшуюся долю:
\[
30\%=\frac{30}{100}=\frac3{10},\qquad
70\%=\frac{70}{100}=\frac7{10}.
\]
Следовательно, время старого комбайна равно
\[
63:\frac7{10}=63\cdot\frac{10}{7}=90\ \text{часов}.
\]

\tagline{Полное правильное решение}
\begin{enumerate}[leftmargin=*,itemsep=2pt]
\item Принимаем время старого комбайна за 100\%.
\item Новый комбайн тратит \(100\%-30\%=70\%\) этого времени.
\item Записываем: \(70\%\) соответствуют 63 часам.
\item Находим 1\%: \(63:70=0{,}9\) часа.
\item Находим 100\%: \(0{,}9\cdot100=90\) часов.
\end{enumerate}

\tagline{Проверка}
Тридцать процентов от 90 часов составляют \(90\cdot0{,}3=27\) часов. Новый комбайн работает на 27 часов быстрее: \(90-27=63\) часа. Условие выполнено.

\finaltag\ Старому комбайну требовалось \(\boxed{90}\) часов.

\tagline{Задача на закрепление}
Новый принтер выполняет работу за 48 минут, что на 20\% меньше времени, необходимого старому принтеру. Сколько минут требовалось старому принтеру?

\end{document}